% Complete Section 7.5, original PDF 98-100 / printed 88-90.
% block: finite-element-space-properties-heading page:98 type:section
\subsection{\BT{Properties of the finite element spaces}{有限元空间的性质}}\label{sec:finite-element-space-properties}

% block: finite-element-approximation-introduction page:98 type:ordinary-paragraph
\BT{The approximation properties of the finite element spaces $\mathcal P_r\Lambda^k(\mathcal T_h)$ and $\mathcal P_r^-\Lambda^k(\mathcal T_h)$ (asymptotically, as the mesh size $h$ tends to $0$) can be established as is usual for finite elements, using the fact that the canonical projection preserves polynomials locally, and applying the Bramble–Hilbert lemma and scaling. This is carried out in detail in [\BookNumber{11}, Theorem 5.3], where the following result is proved.}{有限元空间 $\mathcal P_r\Lambda^k(\mathcal T_h)$ 和 $\mathcal P_r^-\Lambda^k(\mathcal T_h)$ 的逼近性质 (当网格尺寸 $h$ 趋于 $0$ 时的渐近性质), 可像通常的有限元分析那样, 利用典范投影局部保持多项式这一事实, 应用 Bramble–Hilbert 引理和尺度变换来建立. [\BookNumber{11}, Theorem 5.3] 详细进行了这一分析, 并证明了下述结果.}

% block: canonical-projection-approximation-theorem page:98 type:theorem
\Needspace{8\baselineskip}
\begin{Theorem}\label{thm:canonical-projection-approximation}
\BT{Denote by $\Pi_h$ the canonical projection of $\Lambda^k(\Omega)$ onto either $\mathcal P_r\Lambda^k(\mathcal T_h)$ or $\mathcal P_{r+1}^-\Lambda^k(\mathcal T_h)$. Let $1\leq p\leq\infty$ and $(n-k)/p<s\leq r+1$. Then $\Pi_h$ extends boundedly to $W_p^s\Lambda^k(\Omega)$, and there exists a constant $C$ independent of $h$, such that}{以 $\Pi_h$ 表示从 $\Lambda^k(\Omega)$ 到 $\mathcal P_r\Lambda^k(\mathcal T_h)$ 或 $\mathcal P_{r+1}^-\Lambda^k(\mathcal T_h)$ 的典范投影. 设 $1\leq p\leq\infty$ 且 $(n-k)/p<s\leq r+1$. 则 $\Pi_h$ 可有界地延拓至 $W_p^s\Lambda^k(\Omega)$, 且存在不依赖于 $h$ 的常数 $C$, 使得}
\[
\|\omega-\Pi_h\omega\|_{L_p\Lambda^k(\Omega)}\leq Ch^s|\omega|_{W_p^s\Lambda^k(\Omega)},\qquad\omega\in W_p^s\Lambda^k(\Omega).
\]
\end{Theorem}

% block: canonical-projection-trace-restriction page:98 type:ordinary-paragraph
\BT{Note that the restriction $s>(n-k)/p$ comes from the Sobolev embedding theorem: it is what is required for functions in $W_p^s$ to having traces on $k$-dimensional faces.}{注意, 限制 $s>(n-k)/p$ 来自 Sobolev 嵌入定理: 这是 $W_p^s$ 中的函数在 $k$ 维面上具有迹所需的条件.}

% block: complete-and-trimmed-form-errors pages:98,99 type:ordinary-paragraph
\BT{For a smooth $k$-form $\omega$ and the complete polynomial space $\mathcal P_r\Lambda^k(\mathcal T_h)$, we may take $p=2$, $s=r+1$ to obtain an $L^2$ estimate of order $r+1$:}{对光滑 $k$-形式 $\omega$ 和完全多项式空间 $\mathcal P_r\Lambda^k(\mathcal T_h)$, 可取 $p=2$, $s=r+1$, 得到 $r+1$ 阶的 $L^2$ 估计:}
\[
\|\omega-\Pi_{r,h}^k\omega\|_{L^2\Lambda^k}\leq Ch^{r+1}|\omega|_{H^{r+1}\Lambda^k}.
\]
\BT{For the trimmed polynomial space $\mathcal P_r^-\Lambda^k(\mathcal T_h)$, which contains the full polynomial space $\mathcal P_{r-1}\Lambda^k(\mathcal T_h)$ but not $\mathcal P_r\Lambda^k(\mathcal T_h)$, the order is reduced to $r$:}{对修剪多项式空间 $\mathcal P_r^-\Lambda^k(\mathcal T_h)$, 它包含完全多项式空间 $\mathcal P_{r-1}\Lambda^k(\mathcal T_h)$, 但不包含 $\mathcal P_r\Lambda^k(\mathcal T_h)$, 阶数降为 $r$:}
\[
\|\omega-\Pi_{r,h}^{-,k}\omega\|_{L^2\Lambda^k}\leq Ch^r|\omega|_{H^r\Lambda^k}.
\]

% block: complete-form-derivative-error page:99 type:ordinary-paragraph
\BT{For the complete polynomial space, we have the cochain relation}{对完全多项式空间, 有上链关系}
\[
d\Pi_{r,h}^k\omega=\Pi_{r-1,h}^{k+1}d\omega,
\]
\BT{and so}{因而}
\[
\|d(\omega-\Pi_{r,h}^k\omega)\|_{L^2\Lambda^{k+1}}=\|d\omega-\Pi_{r-1,h}^{k+1}d\omega\|_{L^2\Lambda^{k+1}}\leq Ch^r|d\omega|_{H^r\Lambda^{k+1}},
\]
\BT{a loss of one order of convergence for $d\omega$ as compared to $\omega$. For the trimmed space the cochain relation is}{相对于 $\omega$, $d\omega$ 的收敛阶损失了一阶. 对修剪空间, 上链关系为}
\[
d\Pi_{r,h}^{-,k}\omega=\Pi_{r,h}^{-,k+1}d\omega,
\]
\BT{so there is no loss of order:}{因此不损失阶数:}
\[
\|d(\omega-\Pi_{r,h}^{-,k}\omega)\|_{L^2\Lambda^{k+1}}\leq Ch^r|d\omega|_{H^r\Lambda^{k+1}}.
\]

% block: four-finite-element-pairs page:99 type:ordinary-paragraph
\BT{We are now ready to choose the pair of Galerkin subspaces $V_h^{k-1}$, $V_h^k$. Two obvious choices are}{现在可以选择一对 Galerkin 子空间 $V_h^{k-1}$, $V_h^k$. 两种显然的选择是}
\[
\mathcal P_r\Lambda^{k-1}(\mathcal T_h),\mathcal P_{r-1}\Lambda^k(\mathcal T_h)\qquad\text{and}\qquad\mathcal P_r^-\Lambda^{k-1}(\mathcal T_h),\mathcal P_r^-\Lambda^k(\mathcal T_h),
\]
\BT{using either complete polynomial spaces for both subspaces or trimmed spaces for both. There are two other choices as well, which use one complete space and one trimmed. We may reduce $V_h^{k-1}$ in the complete polynomial pair to obtain}{即两个子空间均用完全多项式空间, 或均用修剪空间. 还有另外两种选择, 分别使用一个完全空间和一个修剪空间. 可将完全多项式空间对中的 $V_h^{k-1}$ 缩小, 得到}
\[
\mathcal P_r^-\Lambda^{k-1}(\mathcal T_h),\mathcal P_{r-1}\Lambda^k(\mathcal T_h).
\]
\BT{This does not change the range $dV_h^{k-1}$ and again the canonical projections commute: $d\Pi_{r,h}^{-,k}=\Pi_{r-1,h}^kd$. Similarly, we may expand $V_h^{k-1}$ in the trimmed polynomial pair to get}{这不改变像 $dV_h^{k-1}$, 且典范投影仍然交换: $d\Pi_{r,h}^{-,k}=\Pi_{r-1,h}^kd$. 类似地, 可将修剪多项式空间对中的 $V_h^{k-1}$ 扩大, 得到}
\[
\mathcal P_r\Lambda^{k-1}(\mathcal T_h),\mathcal P_r^-\Lambda^k(\mathcal T_h),
\]
\BT{which, again, does not affect the range of $d$. Altogether, we obtain four possible choices for each degree $r$:}{这同样不影响 $d$ 的像. 总之, 对每个次数 $r$, 得到四种可能的选择:}
\begin{equation}\label{eq:four-galerkin-form-pairs}
V_h^{k-1}=\left\{\begin{gathered}\mathcal P_r\Lambda^{k-1}(\mathcal T_h)\\\text{or}\\\mathcal P_r^-\Lambda^{k-1}(\mathcal T_h)\end{gathered}\right\},\qquad
V_h^k=\left\{\begin{gathered}\mathcal P_r^-\Lambda^k(\mathcal T_h)\\\text{or}\\\mathcal P_{r-1}\Lambda^k(\mathcal T_h)\ (\text{if }r>1)\end{gathered}\right\}.
\end{equation}

% block: galerkin-form-pair-error-rates page:99 type:ordinary-paragraph
\BT{The $L^2$ rates of approximation they achieve for the variables $u\in\Lambda^k$ and $\sigma=\delta u\in\Lambda^{k-1}$ are shown in Table~\ref{tab:galerkin-form-approximation-rates}. That these estimates hold not only for the projection of $\sigma$ and $u$ but also for the Galerkin solution (when the true solution is smooth) can be deduced from the improved error estimates of Theorem~\ref{thm:mixed-improved-errors}.}{这些选择对变量 $u\in\Lambda^k$ 和 $\sigma=\delta u\in\Lambda^{k-1}$ 所达到的 $L^2$ 逼近阶见 Table~\ref{tab:galerkin-form-approximation-rates}. 由 Theorem~\ref{thm:mixed-improved-errors} 的改进误差估计, 可推出这些估计不仅对 $\sigma$ 和 $u$ 的投影成立, 而且对 Galerkin 解也成立 (当真解光滑时).}

% block: third-galerkin-form-space page:99 type:ordinary-paragraph
\BT{Of course, to obtain these results we need to fulfill the hypotheses of Theorem~\ref{thm:mixed-improved-errors}. For this we have to specify a third space $V_h^{k+1}\subset H\Lambda^{k+1}$. (This space is not used in the Galerkin method but is required for theory.) For this we take}{当然, 要得到这些结果, 需要满足 Theorem~\ref{thm:mixed-improved-errors} 的假设. 为此必须指定第三个空间 $V_h^{k+1}\subset H\Lambda^{k+1}$. (Galerkin 方法不使用这个空间, 但理论需要它.) 为此取}
\[
V_h^{k+1}=\begin{cases}\mathcal P_r^-\Lambda^{k+1}(\mathcal T_h)&\text{if }V_h^k=\mathcal P_r^-\Lambda^k(\mathcal T_h),\\\mathcal P_{r-1}^-\Lambda^{k+1}(\mathcal T_h)&\text{if }V_h^k=\mathcal P_{r-1}\Lambda^k(\mathcal T_h).
\end{cases}
\]

% block: galerkin-form-approximation-rates-table page:100 type:table
\begin{table}[htbp]\centering
\caption{\BT{Rates of approximation for stable choices of Galerkin subspaces.}{稳定的 Galerkin 子空间选择所对应的逼近阶.}}\label{tab:galerkin-form-approximation-rates}
\begin{tabular}{cc|cccc}\hline
$V_h^{k-1}$&$V_h^k$&$\sigma$&$d\sigma$&$u$&$du$\\\hline
$\mathcal P_r\Lambda^{k-1}$&$\mathcal P_r^-\Lambda^k$&$r+1$&$r$&$r$&$r$\\
$\mathcal P_r^-\Lambda^{k-1}$&$\mathcal P_r^-\Lambda^k$&$r$&$r$&$r$&$r$\\
$\mathcal P_r\Lambda^{k-1}$&$\mathcal P_{r-1}\Lambda^k$&$r+1$&$r$&$r$&$r-1$\\
$\mathcal P_r^-\Lambda^{k-1}$&$\mathcal P_{r-1}\Lambda^k$&$r$&$r$&$r$&$r-1$\\\hline
\end{tabular}
\end{table}

% block: bounded-form-cochain-projections page:100 type:ordinary-paragraph
\BT{This defines the subcomplex \eqref{eq:galerkin-subcomplex} of the corresponding segment of the de Rham complex, required by Theorem~\ref{thm:mixed-improved-errors}. The remaining hypothesis is that there exist $W$-bounded cochain projections, i.e., bounded projection operators $\pi_h^j:L^2\Lambda^j\to V_h^j$ for $j=k-1$, $k$, $k+1$ which commute with $d$ on its domain. We have already remarked that the canonical projections, defined through the DOFs, form a cochain projection, that is, they commute with $d$ on smooth enough forms. However, the canonical projection is not in general bounded on $L^2\Lambda^k$, nor even on $H\Lambda^k$. For example, $\Pi_{1,h}^0$ is the usual interpolant into the space of piecewise linear functions. It is not bounded on $H\Lambda^0=H^1$ in more than one dimension, because it depends on the point values at the vertices, and it is not $L^2$ bounded even in one dimension. However, there are several ways a bounded cochain projection can be constructed. One approach, which is presented in detail in [\BookNumber{11}, Section 5.4], is to apply a mollifier before applying the canonical projection (following \BookCite{28}), and then to correct the result to obtain a projection with the desired properties (following \BookCite{66}). This construction provides the desired $W$-bounded cochain projections, and hence the error estimates of Theorems~\ref{thm:mixed-basic-error} and~\ref{thm:mixed-improved-errors} hold.}{这就定义了 Theorem~\ref{thm:mixed-improved-errors} 所要求的, 相应 de Rham 复形片段的子复形 \eqref{eq:galerkin-subcomplex}. 剩下的假设是存在 $W$-有界上链投影, 即对 $j=k-1$, $k$, $k+1$, 存在有界投影算子 $\pi_h^j:L^2\Lambda^j\to V_h^j$, 在其定义域上与 $d$ 交换. 前面已经指出, 由自由度定义的典范投影构成上链投影, 即在充分光滑的形式上与 $d$ 交换. 然而, 典范投影一般在 $L^2\Lambda^k$ 上无界, 甚至在 $H\Lambda^k$ 上也无界. 例如, $\Pi_{1,h}^0$ 是到分片线性函数空间的通常插值算子. 在超过一维时, 它在 $H\Lambda^0=H^1$ 上无界, 因为它依赖顶点处的点值; 即使在一维时, 它也不是 $L^2$ 有界的. 不过, 有多种方法可构造有界上链投影. [\BookNumber{11}, Section 5.4] 详细介绍的一种方法, 是先施加磨光算子, 再施加典范投影 (遵循 \BookCite{28}), 然后修正所得结果, 得到具有所需性质的投影 (遵循 \BookCite{66}). 这一构造给出了所需的 $W$-有界上链投影, 因而 Theorem~\ref{thm:mixed-basic-error} 和 Theorem~\ref{thm:mixed-improved-errors} 的误差估计成立.}

% block: cochain-projection-further-properties page:100 type:ordinary-paragraph
\BT{We remark that there are various other properties one might desire from the cochain projections which are not provided by the construction just recalled. For example, one might want the construction to be local or one might ask that the bounds be uniform in the degree $r$. For this reason, the construction of bounded cochain projections remains an active research area \BookCite{30}, \BookCite{43}.}{需要指出, 对上链投影还可能希望有其他各种性质, 而刚才回顾的构造并不提供这些性质. 例如, 可能希望该构造是局部的, 或要求界关于次数 $r$ 一致. 因此, 有界上链投影的构造仍是一个活跃的研究领域 \BookCite{30}, \BookCite{43}.}
